\cotefile{11}
\scanpage{43}
\margin{42}
\mspage{-3.7-}

(identifying $\mathbb{A}_R$ as $\mathbb{A}^{n\nu-n}$), so that $(\mathbb{A}^{n\nu-n},\phi_R)$ is a geometric quotient of $U_R$ by $\mathrm{PGL}(n)$. Moreover, recall that
\[
\sigma_R:\mathbb{A}^{n\nu-n}\longrightarrow U_R
\]
is a global section for $\phi_R$. The following is obvious:

\uline{Corollary 3.2}: Suppose $U\subset U_R$ is any invariant open subset. Then $U\cong\mathrm{PGL}(n)\times\sigma_R^{-1}(U)$, the isomorphism being compatible with the obvious operations of $\mathrm{PGL}(n)$.

To illustrate these results, which may be somewhat obscured by the cumbersome machinery of $R$-partitions, look at the case $\nu=1$, $S_1=\{0,1,\ldots,n\}$. The proposition is then the classical result that $n+2$ points in ``general'' position in $\mathbb{P}_n$ can be normalized to the canonical points $(1,0,\ldots,0)$, $(0,1,\ldots,0),\ldots,(0,0,\ldots,1)$, and $(1,1,\ldots,1)$ by a unique projective transformation. To give a second example, suppose $n=2$, $\nu=2$,
\[
S_1=\{0,1\},\quad S_2=\{1,2\}.
\]
The proposition asserts that, given 5 points $P_0,P_1,P_2$, $P_3$ and $P_4$ in the plane, suppose $P_0,P_1$ and $P_2$ are not collinear, that $P_3$ is not on the line $\overline{P_0P_2}$ or the line $\overline{P_1P_2}$ and that $P_4$ is not on the line $\overline{P_0P_1}$ or $\overline{P_0P_2}$: then we can uniquely normalize $P_0,P_1,P_2$ to $(1,0,0)$, $(0,1,0)$ and $(0,0,1)$ respectively, and normalize $P_3$ to $(1,1,a)$, $P_4$ to $(b,1,1)$ by a suitable projective transformation. (cf. Figure 3, where $Q_3=(1,1,0)$, $Q_4=(0,1,1)$).
